\section{The Channels}
\label{sec:channels}

Table~\ref{tab:channel-map} records the eight named mappings and, importantly, their
different roles and evidence. Four carry \CLEWS{} outputs into \OGCore{}, two apply
scenario policy inputs, and two write \OGCore{} outputs back to \CLEWS{}. The latter are not yet
feedback mechanisms: they emit inputs for a subsequent resource-model solve. The
subsections below concentrate on the economic interpretation and the non-obvious
accounting choices rather than repeat the interface definitions.

\begin{table}[t]
\centering\footnotesize
\caption{The eight interface mappings and their current evidence, grouped by direction
or policy role.}
\label{tab:channel-map}
\begin{tabularx}{\linewidth}{@{}>{\bfseries\raggedright\arraybackslash}p{0.14\linewidth} >{\raggedright\arraybackslash}p{0.18\linewidth} >{\raggedright\arraybackslash}p{0.19\linewidth} >{\raggedright\arraybackslash}p{0.18\linewidth} L@{}}
\toprule
Mapping & Interface quantity & Receiving-model entry & Economic role & Current evidence \\
\midrule
\multicolumn{5}{@{}l}{\emph{CLEWS to OG-Core}} \\
Electricity cost & reform/base cost ratio & household $\tauc$; industry $\TFP_m$ & final- and intermediate-use exposure & composite coupled run; household route isolated \\
Public investment & grid and transmission capex & $\alphaI\to\Kg$ & productive public capital and financing & coupled run; Philippine delta near zero \\
Capital intensity & fleet capital-cost share & energy-industry $\gamma_E$ & production factor shares & aggregate golden record; earlier sectoral identity \\
Health & PM\textsubscript{2.5} path & $\rho(s,t)$; $e(t,s,j)$ & survival and effective labour & controlled runs and coupled run \\
\addlinespace
\multicolumn{5}{@{}l}{\emph{Scenario policy inputs}} \\
Carbon price & exogenous price path & \CLEWS{} penalty; conditional household wedge & emissions cost and revenue & macro run and emitted penalty; no policy re-solve \\
Energy capex & investment-tax-credit path & energy-industry user cost & private generation investment & mapping defined; not exercised here \\
\addlinespace
\multicolumn{5}{@{}l}{\emph{OG-Core to CLEWS}} \\
Discount rate & equilibrium return & \texttt{DiscountRate} & planning-rate harmonisation & artifact emitted; not returned through loop \\
Demand & $Y_m$ or $C_i$ ratio & annual demand path & activity-consistent service demand & re-solve seam tested with exogenous patch \\
\bottomrule
\end{tabularx}
\end{table}

\subsection{Energy price}\label{sec:ch-energy}
The energy-price channel is the composite of \Cref{sec:routes}: it carries the \CLEWS{}
electricity-cost proxy into the economy through both of electricity's uses---a cost-push on
the industries that buy it as an input (the contractionary, wage-bearing side) and a
recycled consumption-tax wedge on households (the regressive, cost-of-living side), with
electricity's self-use netted out so the two do not overlap. The interface uses a
reform/base electricity-cost ratio from the ``auto'' source: a curated cost index where
one exists, otherwise LCOE reconstructed from the run's costs and generation. The
commodity-balance shadow price is degenerate here and is not used. This is a
scenario-derived average-cost proxy rather than an illustrative shock or an observed
retail price. For the Philippine reform the proxy rises about $5\%$ on average.
Incidence is the channel's second output: with
energy a partial necessity ($\cmin>0$), the recycled price change (the
\chan{clean\_incidence} variant) is regressive---the lowest-income group bears more of
it. The standalone \chan{energy\_price} function applies only the household wedge
(\Cref{eq:routeA}), used to isolate the demand response and incidence; the controlled
$+20\%$ stress runs serve the transmission comparison.

\subsection{Carbon price}\label{sec:ch-carbon}
The policy interface can write an exogenous carbon-price path to \CLEWS{} as an
emissions penalty and, where the affected household expenditure is not already priced
through the transmitted electricity-cost ratio, to \OGCore{} as an add-on to the
energy-good wedge,
$\Delta\tauc_{i_e}=\text{(carbon price)}\times\text{deflator}\times\text{carbon
intensity}/p^{b}_{i_e}$, and in \CLEWS{} as an emissions-penalty path. Because \OGCore{}
has no energy in production, only \emph{household} energy carbon is priced on the macro
side; industrial carbon is priced only in \CLEWS{}. If the carbon penalty is already
embodied in the \CLEWS{} cost ratio, the macroeconomic wedge is not added to the same
electricity expenditure (\Cref{sec:discipline}). In the present one-pass demonstration
the \CLEWS{}-side path is \emph{emitted} as a ready input; it does not shape the reported
build mix because the policy path has not been re-solved. Revenue is recyclable through
transfers. Magnitudes are illustrative until the unit/deflator bridge and the carbon
intensity are calibrated (\Cref{sec:calibration}).

\subsection{Investment (public capital)}\label{sec:ch-investment}
Public grid and transmission capex enters as an increment to the public-investment
share $\alphaI$ (as a fraction of GDP, on a finite ramp), accumulating into the public
capital stock $\Kg$ through its law of motion; $\Kg$ is productive via the public-capital
share $\gamma_{g}$. This is the channel in which crowding-out legitimately lives: if the
spending is debt-financed it competes for the savings pool, raising $r$ and displacing
private investment. In the Philippine scenario the public grid-capex delta is near zero
(the transition is generation-side), so the channel is close to inert here. The coupled
run confirms that the mapping is accepted by the economy solve; the Philippine case does
not provide a material public-investment response against which to validate its magnitude.

\subsection{Capital intensity vs.\ capital demand: a worked example in channel discipline}
\label{sec:ch-capint}

Two channels represent the capital side of a capex-heavy generation build-out, and
separating them cleanly is instructive. The temptation is to read ``a clean power system
is capital-intensive'' as ``it pulls capital in.'' \OGCore{}'s own equilibrium conditions
show these are two structurally different instruments---the natural lever for the first
reading need not deliver the second---and give the framework a way to \emph{attribute},
rather than assert, the capital response.

\paragraph{The factor-share lever (\chan{capital\_intensity}).}
This channel raises the energy industry's capital
exponent $\gammam$ by the \CLEWS{} reform/base ratio of the power fleet's capital-cost
share. Under the calibration's Cobb--Douglas technology ($\varepsilon=1$), capital's
revenue share \emph{is} $\gammam$, so the firm's capital demand is
\begin{equation}
  K_m \;=\; \frac{(1-\taub_{m})\,\gammam\, p_m\, Y_m}
                 {\,r+\delta-\taub_{m}\delta^{\tau}_{m}-\tau^{\mathrm{inv}}_{m}\delta\,},
  \label{eq:Kfoc}
\end{equation}
an exact rearrangement of the capital first-order condition. In the controlled comparison,
the channel scales $\gammam$ while the user-cost terms and $(1-\taub)$ factor are unchanged.
Let hats denote reform/base gross ratios. Their cancellation leaves the identity
\begin{equation}
  \hat K_E \;=\; \hat\gamma_E \,\cdot\, \hat p_E \,\cdot\, \hat Y_E .
  \label{eq:Khat}
\end{equation}
The identity is the tool for \emph{reading} the result. On an earlier M=8 Philippine
steady-state solve of this lever
(\Cref{tab:capital-pair}), raising the energy capital exponent by $12.2\%$ turns out to
\emph{shed} energy capital, by $22.4\%$: the higher exponent cuts the sector's unit cost,
so under zero profit the electricity price collapses $42.8\%$; demand is elastic enough
that energy output \emph{rises} $20.9\%$; but the price fall shrinks the sector's revenue
$p_m Y_m$ by more than the exponent and output together add back, so capital---its
$\gammam$ share of that revenue---falls. The identity
$\hat K_E=\hat\gamma_E\hat p_E\hat Y_E$ reproduces the $-22.4\%$ to within $0.03\%$. The
sign is thus not obvious a priori---raising the capital \emph{share} does not draw capital
in---and the framework attributes the response exactly. \chan{capital\_intensity} is a
factor-share and energy-price lever, not a capital-draw-in.

\paragraph{The capital-demand lever (\chan{energy\_capex}).}
The capital-draw-in force is a structurally different instrument
(a defined channel, not exercised as a standalone run in
the present battery): an investment tax credit lowers the energy sector's \emph{cost of
capital}---it enters the user-cost denominator of \cref{eq:Kfoc} through
$\tau^{\mathrm{inv}}$, where $\gammam$ is \emph{absent}---shifting capital demand outward
at the prevailing $r$. That $\gammam$ moves factor \emph{demand} while the ITC moves the
\emph{cost of capital} is a structural fact of the firm block (\Cref{sec:og}), independent
of calibration: they are genuinely different levers, and the no-double-counting principle
(\Cref{sec:discipline}) forbids treating them as substitutes.

\begin{table}[t]
\centering\footnotesize
\caption{The \chan{capital\_intensity} channel decomposed via the firm identity
\cref{eq:Khat}, measured on an earlier M=8 Philippine steady-state solve of the lever
(energy industry $m_e=2$). Raising
the capital exponent collapses the electricity price and, despite higher output, sheds
capital; the identity reproduces the observed capital response to $0.03\%$. The committed
golden records aggregates only (on the re-blessed record the standalone run moves energy
output $+4.4\%$ and aggregate output $+0.001\%$); the sectoral price/capital
decomposition is re-derived per solve.}
\label{tab:capital-pair}
\begin{tabular}{@{}lrrrrr@{}}
\toprule
 & $\hat\gamma_E$ & $\hat p_E$ & $\hat Y_E$ & $\hat K_E$ (obs.) & identity $\hat\gamma_E\hat p_E\hat Y_E$ \\
\midrule
energy industry & $+12.2\%$ & $-42.8\%$ & $+20.9\%$ & $-22.4\%$ & $-22.4\%$ \\
\bottomrule
\end{tabular}
\end{table}

\paragraph{The lesson.}
The two levers act on different objects---a production exponent and the user cost of
capital---and may coexist only when calibrated to distinct technology and policy inputs.
They should not be treated as alternative representations of the same observed capex
change. The firm's first-order condition, \cref{eq:Kfoc}, provides the relevant test: a
higher capital share does not by itself imply that the sector draws in more capital.

\subsection{Health: emissions to demography and labour}\label{sec:ch-health}
The \chan{health} channel maps the reform/base PM\textsubscript{2.5} emissions ratio,
through external Global Burden of Disease age profiles, onto two demographic objects.
\emph{Mortality}: a shock to age-specific mortality $\rho(s,t)$, with the scale calibrated
to hit a signed excess-deaths target (a cleaner reform implies a negative target---lives
saved), after which the population is recomputed. Because PM\textsubscript{2.5} mortality
skews elderly, saved lives are mostly retirees, so the long-run output effect is near
zero. \emph{Morbidity}: the effective-labour profile $e(t,s,j)$ is scaled by a
working-age disability shape, so a cleaner reform raises labour productivity---the main
macro gain. The sign of the GDP effect therefore depends on horizon (a near-term morbidity
gain, a steady-state composition that adds retirees). The Philippine
ambient-PM\textsubscript{2.5} burden and its age profile come from the Global Burden of
Disease export ($43{,}951$ attributable deaths). Translating the power-sector emissions
ratio into that burden uses $M\approx0.082$, constructed from an energy mass share and a
concentration-response elasticity. The lives effect is
$\Delta\text{deaths}=\text{total}\times M\times\Delta\text{emissions}$, so the reform's
$-2.7\%$ emissions change (the first-decade mean) saves about $95$ lives---the
dose-response $M$ is what
separates this from the naive $43{,}951\times\Delta\text{emissions}\approx-1{,}165$. What
remains uncertain---and keeps the channel reduced-form---is the construction of $M$, the
linearity of the emissions-to-burden mapping, and the morbidity productivity elasticity.
The burden and age profile are sourced data; the subsequent translation is a calibrated
model representation. The resulting demographic mechanism requires an age-structured
equilibrium model.

\subsection{Discount rate (economy to resource model)}\label{sec:ch-discount}
After the reform solves, this channel reads \OGCore{}'s equilibrium market return $\rp$
(a multi-period mean) and writes it to \CLEWS{}'s \texttt{DiscountRate}, so that energy
planning uses the model-consistent cost of capital: a lower $\rp$ favours capital-heavy,
long-lived assets. The mapping is a \emph{harmonisation convention}, not a structural
equality: \texttt{DiscountRate} is the rate at which the least-cost planner discounts
system cost, while $\rp$ is the economy's equilibrium portfolio return, and equating
them asserts that energy planning should discount at the economy's opportunity cost of
capital. It does not identify a normative social discount rate. Risk premia, tax
treatment, term structure, and real/nominal conventions remain explicit harmonisation
choices rather than resolved economics. The effect on the macro
economy returns only once the soft link
iterates (\Cref{sec:architecture}); standalone, it is inert by construction.

\subsection{Demand (economy to resource model)}\label{sec:ch-demand}
Symmetrically, this channel reads a reform/base activity ratio (industry output $Y_m$ or
consumption $C_i$, raised to an elasticity) and scales \CLEWS{}'s annual demand. In a
single pass the ratio is $\approx1$, so the channel is latent; it becomes a genuine driver
only inside the iterated loop---the macro-to-energy quantity link in the
\textsc{message-macro}/\textsc{times-macro} tradition
\citep{messner_schrattenholzer_2000,loulou_times_2008}.

\subsection{Policy inputs}\label{sec:levers}
The policy interface uses existing model parameters rather than creating new economic
objects. The investment incentive behind \chan{energy\_capex} sets the relevant
investment-tax-credit and business-tax parameters; the capital-intensity mapping changes
$\gammam$. As \Cref{sec:ch-capint} shows, these parameters act on different equations and
must be tied to distinct scenario quantities if both are used.
